\documentclass[a4paper,landscape,12pt]{extarticle}
\usepackage[left=0.5in,right=0.5in,top=0.4in,bottom=0.55in]{geometry}

\usepackage{graphicx}
\usepackage[normalem]{ulem}
\usepackage{multicol}
\usepackage{array}
\usepackage{caption}
\usepackage{fancyhdr}
\usepackage{lastpage}
\usepackage{amsmath,amssymb}
\usepackage{booktabs}
\usepackage{listings}
\usepackage{xcolor}
\usepackage{hyperref}

% ---------- Listings (code blocks) ----------
\lstset{
  language=Python,
  basicstyle=\ttfamily\Large,
  keywordstyle=\color{blue}\bfseries,
  commentstyle=\color{gray},
  stringstyle=\color{teal},
  numbers=left,
  numberstyle=\Large\color{gray},
  frame=single,
  breaklines=true,
  breakatwhitespace=true,
  columns=fullflexible,
  keepspaces=true,
  showstringspaces=false
}

% Page style with slide numbers
\pagestyle{fancy}
\fancyhf{}
\renewcommand{\headrulewidth}{0pt}
\fancyfoot[R]{\Large\thepage\ / \pageref*{LastPage}}

% Header (expects cpp_logo.png in same folder)
\newcommand{\slideheader}[2]{%
\noindent
\includegraphics[height=0.5in]{cpp_logo.png}%
\hspace{0.3in}%
{\Large \uline{#1 \hfill #2}}%
\par\vspace{0.4in}
}

\begin{document}
\begin{huge}

%%%%%%%%%%%%%%%%%%%%%%%%%%%%%%%%%%%%%%%%%%%%%%%%%%%%%%%%%%%%
% SLIDE 1 - Title Slide
%%%%%%%%%%%%%%%%%%%%%%%%%%%%%%%%%%%%%%%%%%%%%%%%%%%%%%%%%%%%
\slideheader{ECE 5110 Numerical Modeling}{}

\vspace{1.0in}
\begin{center}
{\Huge \textbf{Solving Linear Systems of Equations}}

\vspace{0.35in}
{\LARGE Unit 4: Gaussian Elimination \& Cubic Splines}\\
\vspace{0.15in}
{\LARGE \textbf{From MATLAB to Python — solve\_lsoe + cubic\_spline}}

\vspace{0.4in}
{\Large Jared Acosta, Josiah Wiggins, Jake Davila, Jack Gaon, Jacky Li, Virgil Nhieu, Elijah Perez, Richie Wong}

\vspace{0.15in}
{\Large April 13, 2026}
\end{center}

\newpage

%%%%%%%%%%%%%%%%%%%%%%%%%%%%%%%%%%%%%%%%%%%%%%%%%%%%%%%%%%%%
% SLIDE 2 - Outline
%%%%%%%%%%%%%%%%%%%%%%%%%%%%%%%%%%%%%%%%%%%%%%%%%%%%%%%%%%%%
\slideheader{Outline}{Linear Systems}

\leftskip 0.7in \rightskip 0.7in
\vspace{0.2in}

\begin{itemize}
  \item Why Solve $\mathbf{A}\mathbf{x} = \mathbf{b}$?
  \item Gaussian Elimination — Four Steps
  \item Back Substitution (the missing piece)
  \item Cubic Spline Interpolation Theory
  \item Building the Spline Constraint Matrix
  \item Implementation in \texttt{tools.py}
  \item Test Cases \& Results
  \item Discussion \& Conclusion
\end{itemize}

\leftskip 0in \rightskip 0in
\newpage

%%%%%%%%%%%%%%%%%%%%%%%%%%%%%%%%%%%%%%%%%%%%%%%%%%%%%%%%%%%%
% SLIDE 3 - Why Solve Ax = b?
%%%%%%%%%%%%%%%%%%%%%%%%%%%%%%%%%%%%%%%%%%%%%%%%%%%%%%%%%%%%
\slideheader{Why Solve $\mathbf{A}\mathbf{x} = \mathbf{b}$?}{Motivation}

\leftskip 0.7in \rightskip 0.7in
\vspace{0.15in}

\begin{multicols}{2}
\textbf{Applications:}
\begin{itemize}
  \item Circuit analysis (KVL/KCL)
  \item Structural engineering (FEM)
  \item Least-squares fitting
  \item Control systems
  \item Graphics \& transformations
\end{itemize}

\columnbreak

\textbf{The Problem:}
\[
\underbrace{\begin{bmatrix} a_{11} & \cdots & a_{1n} \\ \vdots & \ddots & \vdots \\ a_{n1} & \cdots & a_{nn} \end{bmatrix}}_{\mathbf{A}}
\underbrace{\begin{bmatrix} x_1 \\ \vdots \\ x_n \end{bmatrix}}_{\mathbf{x}}
= \underbrace{\begin{bmatrix} b_1 \\ \vdots \\ b_n \end{bmatrix}}_{\mathbf{b}}
\]

\vspace{0.15in}
Find $\mathbf{x}$ given $\mathbf{A}$ and $\mathbf{b}$.
\end{multicols}

\leftskip 0in \rightskip 0in
\newpage

%%%%%%%%%%%%%%%%%%%%%%%%%%%%%%%%%%%%%%%%%%%%%%%%%%%%%%%%%%%%
% SLIDE 4 - Gaussian Elimination Overview
%%%%%%%%%%%%%%%%%%%%%%%%%%%%%%%%%%%%%%%%%%%%%%%%%%%%%%%%%%%%
\slideheader{Gaussian Elimination — Four Steps}{Algorithm}

\leftskip 0.7in \rightskip 0.7in
\vspace{0.15in}

Transform $[\mathbf{A}|\mathbf{b}]$ into $[\mathbf{U}|\mathbf{c}]$ (upper-triangular), then solve by back substitution.

\vspace{0.2in}
\begin{enumerate}
  \item \textbf{Normalize} rows by their max absolute value
  \item \textbf{Pivot} — swap rows to maximize diagonal entry
  \item \textbf{Eliminate} — zero out below-diagonal entries
  \item \textbf{Back-substitute} — solve from bottom up
\end{enumerate}

\vspace{0.2in}
\begin{center}
$\begin{bmatrix} * & * & * \\ * & * & * \\ * & * & * \end{bmatrix}
\xrightarrow{\text{steps 1--3}}
\begin{bmatrix} * & * & * \\ 0 & * & * \\ 0 & 0 & * \end{bmatrix}
\xrightarrow{\text{step 4}}
\begin{bmatrix} x_1 \\ x_2 \\ x_3 \end{bmatrix}$
\end{center}

\leftskip 0in \rightskip 0in
\newpage

%%%%%%%%%%%%%%%%%%%%%%%%%%%%%%%%%%%%%%%%%%%%%%%%%%%%%%%%%%%%
% SLIDE 5 - Steps 1 & 2: Normalize + Pivot
%%%%%%%%%%%%%%%%%%%%%%%%%%%%%%%%%%%%%%%%%%%%%%%%%%%%%%%%%%%%
\slideheader{Step 1: Normalize \quad Step 2: Pivot}{Algorithm}

\leftskip 0.7in \rightskip 0.7in
\vspace{0.1in}

\begin{multicols}{2}
\textbf{Row Normalization:}

Scale each active row so the largest entry has magnitude 1:
\[
A_{r,:} \leftarrow \frac{A_{r,:}}{\max_j |A_{r,j}|}
\]

Prevents large entries from dominating pivot selection.

\columnbreak

\textbf{Partial Pivoting:}

Swap row $k$ with the row below it that has the largest $|A_{r,k}|$:
\[
\text{pivot} = \arg\max_{r \ge k} |A_{r,k}|
\]

Reduces round-off error amplification.
\end{multicols}

\leftskip 0in \rightskip 0in
\newpage

%%%%%%%%%%%%%%%%%%%%%%%%%%%%%%%%%%%%%%%%%%%%%%%%%%%%%%%%%%%%
% SLIDE 6 - Step 3: Forward Elimination
%%%%%%%%%%%%%%%%%%%%%%%%%%%%%%%%%%%%%%%%%%%%%%%%%%%%%%%%%%%%
\slideheader{Step 3: Forward Elimination}{Algorithm}

\leftskip 0.7in \rightskip 0.7in
\vspace{0.1in}

For each row $r > k$, compute the multiplier and subtract:
\[
m_{r,k} = \frac{A_{r,k}}{A_{k,k}}, \qquad
A_{r,:} \leftarrow A_{r,:} - m_{r,k} \cdot A_{k,:}
\]

\vspace{0.15in}
\textbf{Example (instructor's 4×4 system):}
\[
\begin{bmatrix}
1 & 1 & 1 & 1 \\
1 & 2 & 3 & 4 \\
0 & 2 & 0 & 1 \\
-1 & -2 & 1 & 2
\end{bmatrix}
\xrightarrow{\text{eliminate col 1}}
\begin{bmatrix}
1 & 1 & 1 & 1 \\
0 & 1 & 2 & 3 \\
0 & 2 & 0 & 1 \\
0 & -1 & 2 & 3
\end{bmatrix}
\xrightarrow{\cdots}
\mathbf{U}
\]

Repeat for columns $k = 1, 2, \ldots, n{-}1$.

\leftskip 0in \rightskip 0in
\newpage

%%%%%%%%%%%%%%%%%%%%%%%%%%%%%%%%%%%%%%%%%%%%%%%%%%%%%%%%%%%%
% SLIDE 7 - Step 4: Back Substitution (THE MISSING PIECE)
%%%%%%%%%%%%%%%%%%%%%%%%%%%%%%%%%%%%%%%%%%%%%%%%%%%%%%%%%%%%
\slideheader{Step 4: Back Substitution (Missing from MATLAB)}{Algorithm}

\leftskip 0.7in \rightskip 0.7in
\vspace{0.1in}

After forward elimination, we have $\mathbf{U}\mathbf{x} = \mathbf{c}$:

\[
\begin{bmatrix}
U_{11} & U_{12} & \cdots & U_{1n} \\
0 & U_{22} & \cdots & U_{2n} \\
\vdots &  & \ddots & \vdots \\
0 & 0 & \cdots & U_{nn}
\end{bmatrix}
\begin{bmatrix} x_1 \\ x_2 \\ \vdots \\ x_n \end{bmatrix}
= \begin{bmatrix} c_1 \\ c_2 \\ \vdots \\ c_n \end{bmatrix}
\]

\vspace{0.1in}
Solve \textbf{bottom-up}:
\[
\boxed{x_i = \frac{1}{U_{ii}}\left(c_i - \sum_{j=i+1}^{n} U_{ij}\,x_j\right), \quad i = n, n{-}1, \ldots, 1}
\]

\vspace{0.1in}
This was the step \textbf{missing} from the original MATLAB code.

\leftskip 0in \rightskip 0in
\newpage

%%%%%%%%%%%%%%%%%%%%%%%%%%%%%%%%%%%%%%%%%%%%%%%%%%%%%%%%%%%%
% SLIDE 8 - Implementation
%%%%%%%%%%%%%%%%%%%%%%%%%%%%%%%%%%%%%%%%%%%%%%%%%%%%%%%%%%%%
\slideheader{Implementation: solve\_lsoe (tools.py)}{Code}

\leftskip 0.7in \rightskip 0.7in
\vspace{0.05in}

\begin{lstlisting}
def solve_lsoe(self, A, B):
    A = np.array(A, dtype=np.float64)
    B = np.array(B, dtype=np.float64).flatten()
    n = A.shape[0]
    for ix in range(n - 1):
      # 1. Normalize
      for r in range(ix, n):
        A[r,ix:] /= np.max(np.abs(A[r,ix:]))
      # 2. Pivot
      piv = ix + np.argmax(np.abs(A[ix:,ix]))
      A[[ix,piv]] = A[[piv,ix]]
      # 3. Eliminate
      for r in range(ix+1, n):
        m = A[r,ix]/A[ix,ix]
        A[r,ix:] -= m * A[ix,ix:]
    # 4. Back Substitution
    x = np.zeros(n)
    for i in range(n-1, -1, -1):
      x[i] = (B[i] - A[i,i+1:]@x[i+1:]) / A[i,i]
    return x, 0
\end{lstlisting}

\leftskip 0in \rightskip 0in
\newpage

%%%%%%%%%%%%%%%%%%%%%%%%%%%%%%%%%%%%%%%%%%%%%%%%%%%%%%%%%%%%
% SLIDE 9 - Test Cases & Results
%%%%%%%%%%%%%%%%%%%%%%%%%%%%%%%%%%%%%%%%%%%%%%%%%%%%%%%%%%%%
\slideheader{Test Cases \& Results}{Verification}

\leftskip 0.7in \rightskip 0.7in
\vspace{0.1in}

\begin{multicols}{2}
\textbf{Instructor's 4×4 System:}
\[
\mathbf{x} = \begin{bmatrix} 5.5 \\ 1 \\ -12.5 \\ 10 \end{bmatrix}
\]
Max error vs \texttt{numpy}: $1.78 \times 10^{-15}$

\vspace{0.15in}
\textbf{All Tests:}
\begin{center}
\begin{tabular}{lc}
\toprule
\textbf{Test} & \textbf{Status} \\ \midrule
4\texttimes 4 instructor    & $\checkmark$ \\
Identity 3\texttimes 3      & $\checkmark$ \\
Simple 2\texttimes 2        & $\checkmark$ \\
Singular detect   & $\checkmark$ \\
Random 6\texttimes 6        & $\checkmark$ \\
Spline sin (6 knots)  & $\checkmark$ \\
Spline 4-point        & $\checkmark$ \\
\bottomrule
\end{tabular}
\end{center}

\columnbreak

\begin{center}
\includegraphics[width=0.45\textwidth]{unit04_accuracy_analysis.png}
\end{center}

Error stays near machine precision for well-conditioned matrices.
\end{multicols}

\leftskip 0in \rightskip 0in
\newpage

%%%%%%%%%%%%%%%%%%%%%%%%%%%%%%%%%%%%%%%%%%%%%%%%%%%%%%%%%%%%
% SLIDE 10 - Cubic Spline Theory
%%%%%%%%%%%%%%%%%%%%%%%%%%%%%%%%%%%%%%%%%%%%%%%%%%%%%%%%%%%%
\slideheader{Cubic Spline Interpolation}{Theory}

\leftskip 0.7in \rightskip 0.7in
\vspace{0.1in}

Given $n$ data points, fit $n-1$ cubic polynomials:
\[
S_i(x) = a_i x^3 + b_i x^2 + c_i x + d_i
\]

\vspace{0.1in}
\textbf{5 types of constraints} ($4(n{-}1)$ equations total):
\begin{enumerate}
  \item Each spline passes through its \textbf{left} knot: $S_i(x_i) = y_i$
  \item Each spline passes through its \textbf{right} knot: $S_i(x_{i+1}) = y_{i+1}$
  \item $S_i'(x_{i+1}) = S_{i+1}'(x_{i+1})$ — first derivative continuity
  \item $S_i''(x_{i+1}) = S_{i+1}''(x_{i+1})$ — second derivative continuity
  \item Natural BCs: $S''(x_0) = 0$, $S''(x_{n-1}) = 0$
\end{enumerate}

\vspace{0.1in}
Build matrix $\mathbf{A}$ and vector $\mathbf{b}$, solve with \texttt{solve\_lsoe}!

\leftskip 0in \rightskip 0in
\newpage

%%%%%%%%%%%%%%%%%%%%%%%%%%%%%%%%%%%%%%%%%%%%%%%%%%%%%%%%%%%%
% SLIDE 11 - Cubic Spline Code
%%%%%%%%%%%%%%%%%%%%%%%%%%%%%%%%%%%%%%%%%%%%%%%%%%%%%%%%%%%%
\slideheader{cubic\_spline Implementation}{Code}

\leftskip 0.7in \rightskip 0.7in
\vspace{0.05in}

\begin{lstlisting}
def cubic_spline(self, X, Y):
    n = len(X)
    A = np.zeros((4*(n-1), 4*(n-1)))
    B = np.zeros(4*(n-1))
    # 1. Start pts: a*xi^3+b*xi^2+c*xi+d = yi
    # 2. End pts:   same at x_{i+1}
    # 3. 1st deriv: 3a*x^2+2b*x+c match
    # 4. 2nd deriv: 6a*x+2b match
    # 5. Natural: S''=0 at endpoints
    coeffs, err = self.solve_lsoe(A, B)
    return coeffs, err
\end{lstlisting}

Key insight: the spline problem reduces to a \textbf{linear system} solved by our own Gaussian elimination!

\leftskip 0in \rightskip 0in
\newpage

%%%%%%%%%%%%%%%%%%%%%%%%%%%%%%%%%%%%%%%%%%%%%%%%%%%%%%%%%%%%
% SLIDE 12 - Cubic Spline Results
%%%%%%%%%%%%%%%%%%%%%%%%%%%%%%%%%%%%%%%%%%%%%%%%%%%%%%%%%%%%
\slideheader{Cubic Spline Results}{Verification}

\leftskip 0.7in \rightskip 0.7in
\vspace{0.1in}

\begin{center}
\includegraphics[width=0.7\textwidth]{unit04_cubic_spline.png}
\end{center}

\vspace{0.1in}
Cubic spline interpolation of $\sin(x)$ with 6 knots. Error peaks at $\sim 10^{-3}$
between knots — excellent for only 6 data points.

\leftskip 0in \rightskip 0in
\newpage

%%%%%%%%%%%%%%%%%%%%%%%%%%%%%%%%%%%%%%%%%%%%%%%%%%%%%%%%%%%%
% SLIDE 13 - Discussion & Conclusion
%%%%%%%%%%%%%%%%%%%%%%%%%%%%%%%%%%%%%%%%%%%%%%%%%%%%%%%%%%%%
\slideheader{Discussion \& Conclusion}{}

\leftskip 0.7in \rightskip 0.7in
\vspace{0.1in}

\begin{multicols}{2}
\textbf{Key Points:}
\begin{itemize}
  \item Scaled partial pivoting prevents catastrophic cancellation
  \item Overall complexity: $O(n^3)$
  \item Cubic splines build a $4(n{-}1)$ system, solved by \texttt{solve\_lsoe}
  \item Natural BCs: $S'' = 0$ at endpoints
\end{itemize}

\columnbreak

\textbf{What We Did:}
\begin{itemize}
  \item Translated MATLAB $\to$ Python
  \item \textbf{Added back substitution}
  \item \textbf{Added cubic spline} interpolation
  \item 7 test cases, all passing
  \item Accuracy \& spline plots
\end{itemize}
\end{multicols}

\vspace{0.2in}
\begin{center}
\textbf{Conclusion:} Complete linear system solver + cubic spline interpolation,
both verified to machine precision. Spline builds on solve\_lsoe.
\end{center}

\leftskip 0in \rightskip 0in

\end{huge}
\end{document}
